\chapter{Zapiski}

\section{12. 11. 2025}
\begin{verbatim}
Poglej dataset: https://bdsp.io/content/hsp/2.0/
https://www.nature.com/articles/s41467-025-61945-1#data-availability
\end{verbatim}


\section{Prvi poskus}
Vzel sem en dataset (CAP Sleep Database), določil 6 kanalov (F4-C4, C4-P4, P4-O2, F3-C3, C3-P3, P3-O1), ter
na njih filtriral frekvence (raw.filter(l\_freq=0.3, h\_freq=30)). Nato sem vsem spremenil frekvenco vzorčenja na
256 Hz. Podatke sem razdelil na 30 sekundne intervale ter za vsakega shranil spectrogram in Sleep stage (W, S1, S2, S3, S4, REM).
\begin{verbatim}
def get_spectrogram(raw):
    segment_data = raw.get_data()

    spectrograms = []
    for channel_data in segment_data:
        f, t, Zxx = scipy.signal.stft(
            channel_data,
            fs=TARGET_SFREQ,
            nperseg=256,
            noverlap=192
        )
        
        f_mask = (f < 50)
        Zxx = Zxx[f_mask, :]

        spectrogram = np.log1p(np.abs(Zxx))
        spectrograms.append(spectrogram)
    return np.stack(spectrograms, axis=0).astype(np.float32)
\end{verbatim}

Fazi S3 in S4 nato premapiramo v fazo N3.

Če katerega od 6ih kanalov ni, vstavimo vrstico ničel.

Model je bil naslednji:
\begin{verbatim}
class CnnEmbed(nn.Module):
    """Split EEG into fixed-size patches and embed them."""
    def __init__(self, emb_size=64, num_channels=6, height=50, width=122, dropout_prob=0.5):
        super().__init__()
        print(f"Initializing CnnEmbed with emb_size={emb_size}, num_channels={num_channels}, height={height}, width={width}, dropout_prob={dropout_prob}", flush=True)
        self.emb_size = emb_size

        self.cnn_feature_extractor = nn.Sequential(
            nn.Conv2d(num_channels, 16, kernel_size=(3, 3), stride=(2, 2), padding=1),
            nn.ReLU(),
            nn.GroupNorm(4, 16),
            nn.Conv2d(16, emb_size, kernel_size=(3, 3), stride=(2, 2), padding=1),
            nn.ReLU(),
            nn.GroupNorm(8, emb_size),
        )

        with torch.no_grad():
            dummy = torch.zeros(1, num_channels, height, width)
            cnn_out = self.cnn_feature_extractor(dummy)
            flattened_dim = cnn_out.view(1, -1).size(1)

        self.projection = nn.Linear(flattened_dim, emb_size)
        self.dropout = nn.Dropout(p=dropout_prob)

    def forward(self, x):
        cnn_out = self.cnn_feature_extractor(x)
        flattened = cnn_out.view(cnn_out.size(0), -1)
        flattened = self.dropout(flattened)
        token = self.projection(flattened)

        return token

# -----------------------------
# Transformer-based model
# -----------------------------
class EEGTransformer(nn.Module):
    def __init__(self, emb_size=64, nhead=4, num_layers=2, num_classes=5, num_channels=6, height=50, width=122, dropout_prob=0.65):
        super().__init__()
        print(f"Initializing EEGTransformer with emb_size={emb_size}, nhead={nhead}, num_layers={num_layers}, num_classes={num_classes}, num_channels={num_channels}, height={height}, width={width}, dropout_prob={dropout_prob}", flush=True)

        self.emb_size = emb_size

        self.cnn_embed = CnnEmbed(
            emb_size=emb_size,
            num_channels=num_channels,
            height=height,
            dropout_prob=dropout_prob,
            width=width
        )

        self.pos_emb = nn.Embedding(10, emb_size)

        # Standard Transformer Encoder
        encoder_layer = nn.TransformerEncoderLayer(d_model=emb_size, nhead=nhead, batch_first=True)
        self.transformer = nn.TransformerEncoder(encoder_layer, num_layers=num_layers)
        self.attn_pool = nn.Linear(emb_size, 1)

        # Classification head
        self.classification_head = nn.Sequential(
            nn.Linear(emb_size, emb_size),
            nn.ReLU(),
            nn.Dropout(dropout_prob),
            nn.Linear(emb_size, num_classes)
        )

    def forward(self, segments, patient_ids, masks):
        # segments: (batch, 6, 15360)
        batch_size, window_size, num_channels, height, width = segments.shape

        # 1. Patch embedding
        segments_reshaped = segments.view(batch_size * window_size, num_channels, height, width)
        tokens = self.cnn_embed(segments_reshaped)
        tokens = tokens.view(batch_size, window_size, self.emb_size)

        # 2. Add positional embedding
        positions = torch.arange(window_size, device=segments.device)
        pos_emb = self.pos_emb(positions).unsqueeze(0)  # (1, window_size, emb_size)
        tokens = tokens + pos_emb

        # 3. Transformer
        # TODO DODAJ PADDING MASK , src_key_padding_mask=~masks
        transformer_out = self.transformer(tokens)  # (batch, 768, emb_size)

        # 4. Mean + learned pool across tokens
        weights = torch.softmax(self.attn_pool(transformer_out), dim=1)
        pooled_mean = transformer_out.mean(dim=1)
        pooled_attn = (transformer_out * weights).sum(dim=1)
        pooled = 0.5 * pooled_mean + 0.5 * pooled_attn

        # 5. Classify
        logits = self.classification_head(pooled)
        return logits
\end{verbatim}

Torej, vhod je več 30 sekundnih zapisov (najpogosteje window\_size = 10, torej 5 minut),
vsakih 30 sekund gre čez CNN in dobimo 1 token dolžine EMB\_SIZE iz vsakega 30 sekundnega zapisa.
Ti tokeni gredo v Transformer, kjer se upošteva kontekst (prejšnji in naslednji tokeni). (Tokeni imajo dodan še position embedding)

Za stestirat: kakšen vpliv sploh ima dolžina konteksta, izgledal je kot da med windows\_size = 1 in 10 je 
bil profit around 1\% accuracy.

Napovedat poskusimo vsebino srednjega zapisa (npr. če je window size 10, poskusimo napovedat sleep stage za 6ti zapis).

Uporabljal sem focal loss, ker je za N1 zelo malo zapisov.
Najboljše rezultate sem dobil za: 
EMB\_SIZE=32, BATCH\_SIZE=16, LEARNING\_RATE=5e-05, WEIGHT\_DECAY=0.001, EPOCHS=50
Initializing EEGTransformer with emb\_size=32, nhead=4, num\_layers=2, num\_classes=5, num\_channels=6, height=50, width=122, dropout\_prob=0.5

\begin{table}[h!]
\centering
\caption{Confusion Matrix}
\begin{tabular}{c|c|ccccc}
\toprule
 & & \multicolumn{5}{c}{\textbf{Predicted Label}} \\
\cmidrule(lr){3-7}
 & & S & N1 & N2 & N3 & R \\
\midrule
\multirow{5}{*}{\rotatebox{90}{\textbf{Actual Label}}}
 & S   & 2634 & 8   & 185  & 13   & 154 \\
 & N1  & 301  & 0   & 244  & 1    & 213 \\
 & N2  & 134  & 0   & 4656 & 508  & 467 \\
 & N3  & 5    & 0   & 657  & 3164 & 187 \\
 & R   & 54   & 0   & 496  & 23   & 2388 \\
\bottomrule
\end{tabular}
\end{table}

\[
\text{Accuracy} = 0.7787
\]


Sem imel validation set in vzel tisti model, ki se je na validation datasetu najboljše odrezal.
Imel sem tudi early stopping

\begin{verbatim}
criterion = FocalLoss(task_type='multi-class', alpha=None, gamma=2.0, reduction='mean', num_classes=NUM_CLASSES)
optimizer = torch.optim.Adam(model.parameters(), lr=LEARNING_RATE, weight_decay=WEIGHT_DECAY)
scheduler = CosineAnnealingLR(optimizer, T_max=EPOCHS, eta_min=1e-6)
\end{verbatim}

\section{Drugi poskus}

V prvem datasetu sem začel dodajati različne kanale, kadar tisti primarni v podatkih ne obstajajo.
Prav tako sem dodal data augmentation - kanale pred treniranjem premešam, zato da se model ne more naučiti
channel specific značilnosti. To je pripeljalo accuracy 0.831.

\begin{figure}[H]
    \begin{center}
        \includegraphics[width=12cm]{cm/confusion_matrix_cap2_30s_4_trans.png}
    \end{center}
\end{figure}

Nato sem dodal še dva podatkovna seta $\rightarrow$ 
"Adult High-Density \& Wearable EEG" (AHD)
in "Dataset of EEG recordings containing HFO markings for 30 pediatric patients with epilepsy" (HFO)

Model sem naredil malo kompleksenjši, zmanjšal training rate, poskušal prilagoditi parametre večim podatkom.

\begin{verbatim}
Initializing EEGTransformer with emb_size=64, nhead=4, num_layers=4, num_classes=5, num_channels=6, height=50, width=122, dropout_prob=0.65
Initializing CnnEmbed with emb_size=64, num_channels=6, height=50, width=122, dropout_prob=0.65
\end{verbatim}

izkaže se, da je HFO set precej bolj zamašen, večji maximumi, večja varianca. Par razlogov:
\begin{itemize}
    \item otroci imajo močnejšo voltažo pri EEG signalih
    \item za preverit če lahko odstranimo šum kot so mežiki itd.
    \item mogoče tudi podatki niso preveč lepi, nekaterih faz imamo zelo malo, lahko je malo vprašljivo.
\end{itemize}

Za poskusit:
Najprej bom normaliziral inpute kot so, ne pa en skupen mean in std za VSE podatke.
Mogoče bi bilo smiselno imeti mean za otroke in odrasle posebej in normalizirati skupine posebej.
Verjetno se lahko doda še embeddings za starostne skupine
Če najdem še kak podatkovni set z EEG otrok je def. za preverit če so problemi age specific ali zgolj
na tem datasetu HFO.

\begin{figure}[H]
    \begin{center}
        \includegraphics[width=\linewidth]{cm/all_confusion_matrices_1.png}
    \end{center}
\end{figure}

\begin{verbatim}
    Accuracy for source PROCESSED_AHD: 0.7535
    Accuracy for source PROCESSED_CAP_2: 0.8356
    Accuracy for source PROCESSED_HFO: 0.6150
    Joint Accuracy: 0.8004
\end{verbatim}